% Chapter 12: The Isosceles Triangle Case
% ============================================================

\chapter{The Isosceles Triangle Case}
\label{ch:isosceles}

When the triangle is isosceles, the symmetry reduces the system's degrees of freedom and yields a lower-degree polynomial.

\section{Setup}

Without loss of generality, let $b = c$ (the two sides adjacent to vertex $A$ are equal). Then $\tri ABC$ is isosceles with the axis of symmetry passing through $A$ and the midpoint of $BC$.

The Cayley--Menger constraint with $b = c$ becomes:
\begin{equation}
\CM_4\big|_{b=c} = \begin{vmatrix}
0 & 1 & 1 & 1 & 1 \\
1 & 0 & b^2 & b^2 & d_1^2 \\
1 & b^2 & 0 & a^2 & d_2^2 \\
1 & b^2 & a^2 & 0 & d_3^2 \\
1 & d_1^2 & d_2^2 & d_3^2 & 0
\end{vmatrix} = 0.
\end{equation}

\section{Reduction to Two Unknowns}

With $b = c$, we have two unknowns: $a$ and $b$. The Cayley--Menger determinant provides one polynomial equation in $(a, b)$. The isosceles constraint has already consumed one degree of freedom, reducing the deficit from 2 to 1.

\subsection{The Reduced Polynomial}

Expanding $\CM_4|_{b=c}$ (using symbolic computation):

Using the auxiliary variables:
\begin{align}
u &= d_2^2 + d_3^2 - a^2, \\
v &= d_1^2 + d_3^2 - b^2, \\
w &= d_1^2 + d_2^2 - b^2.
\end{align}

Note that $w$ now uses $b$ instead of $c$, since $b = c$.

The master polynomial~\eqref{eq:master-final} becomes:
\begin{equation}
d_3^2 w^2 - uvw + d_2^2 v^2 + d_1^2 u^2 - 4d_1^2 d_2^2 d_3^2 = 0.
\end{equation}

Substituting $w = d_1^2 + d_2^2 - b^2$ and simplifying yields a polynomial in $(a^2, b^2)$ that can be analyzed for fixed $(d_1, d_2, d_3)$.

\section{Symmetry Axis Constraint}

For an isosceles triangle $b = c$, the axis of symmetry passes through vertex $A$ and the midpoint $M$ of $BC$. If $P$ lies on this axis (i.e., $d_2 = d_3$), additional simplification occurs.

\begin{theorem}[Isosceles with Symmetric Point]
If $b = c$ and $d_2 = d_3$, the Cayley--Menger determinant reduces to a quartic in $b^2$ (for fixed $a$) or a quartic in $a^2$ (for fixed $b$).
\end{theorem}

\section{Special Case: $d_2 = d_3$}

When $d_2 = d_3 = d$ (the interior point is equidistant from $B$ and $C$), the problem inherits the bilateral symmetry. Setting $u = 2d^2 - a^2$, $v = d_1^2 + d^2 - b^2$:

\begin{equation}
d^2(d_1^2 + d^2 - b^2)^2 - (2d^2 - a^2)(d_1^2 + d^2 - b^2)^2 + d^2(d_1^2 + d^2 - b^2)^2 + d_1^2(2d^2 - a^2)^2 - 4d_1^2 d^4 = 0.
\end{equation}

This simplifies considerably, yielding a relation between $a$ and $b$ that, combined with $b = c$, determines the triangle up to the remaining degree of freedom.

\section{Perimeter Formula for Isosceles}

The perimeter in the isosceles case ($b = c$) is:
\begin{equation}
\perim = a + 2b.
\end{equation}

Given the polynomial relation $F(a^2, b^2; d_1, d_2, d_3) = 0$, the perimeter becomes a function of one parameter (e.g., $a$), with $b$ determined implicitly by $F = 0$:
\begin{equation}
\perim(a) = a + 2\sqrt{b^2(a)},
\end{equation}
where $b^2(a)$ is obtained by solving $F$ for $b^2$.

\section{Numerical Examples}

\begin{example}
Let $d_1 = 2$, $d_2 = d_3 = 3$, and require $b = c$.

The symmetry $d_2 = d_3$ is consistent with $b = c$ when $P$ lies on the axis of symmetry. The Cayley--Menger constraint gives a polynomial in $(a, b)$ that yields specific solutions.
\end{example}
